\documentclass[12pt,a4paper]{article}
\usepackage{amsmath,amsfonts,amssymb}
\usepackage{graphicx}
\usepackage{hyperref}
\usepackage{enumitem}
\usepackage{geometry}
\geometry{margin=2.5cm}
\emergencystretch=2em

\begin{document}

\begin{center}
    {\Large \textbf{Master Thesis Project in Quantum Information Science and Technology}}\\[0.3cm]
    {\large \textbf{Exploring and Extending the Rodeo Algorithm for Quantum Eigenvalue Problems on Real Quantum Hardware: Multi-Qubit Hamiltonians and Physics-Based Input States}}\\[0.2cm]
    \vspace{0.3cm}
    {Duration: One Year (60 ECTS)}
\end{center}

\vspace{1cm}


\subsection*{Introduction}

Estimating eigenvalues and preparing eigenstates of quantum
Hamiltonians are central tasks in quantum many-body physics, quantum
chemistry, and nuclear structure theory.  The \emph{rodeo algorithm}
(RA)~[1] is a stochastic, interference-based method that suppresses
unwanted spectral components and isolates the eigenstate at a chosen
target energy.  The suppression is exponential in the number of cycles
of the algorithm.  For the preparation of eigenstates this is
exponentially faster than quantum phase estimation and adiabatic
evolution~[1,2], provided that the initial state has a sufficient
overlap with the eigenstate one wants.

The first implementation of the RA on real quantum hardware~[3]
demonstrated high-precision energy estimation for one-qubit
Hamiltonians on an IBM Q device.  The method was validated through
calculations based on the Hellmann--Feynman theorem and through a
comparison with direct preparation of the eigenvector.  A second
demonstration~[4] treated a two-qubit Hamiltonian on the Quantinuum
H1-2 and IBM Perth devices, with controlled reversal gates that
reduced the number of two-qubit gates by a factor of five.  The choice
of the evolution times, which were random in the original formulation,
has since been optimized~[5].

Two things limit the step from these demonstrations to interacting
many-body systems.  The first is the circuit.  Every cycle needs a
controlled time evolution, which for several qubits has to be
Trotterized and executed, together with mid-circuit measurements, on
noisy devices.  The second is the input state.  Like quantum phase
estimation, the RA is a filter: its success probability is
proportional to the overlap of the input state with the wanted
eigenstate.  For simple input states, such as a product state, this
overlap decreases exponentially with the size of the system.  The
rodeo algorithm is therefore only as good as the state it is given.

\emph{Physics-based algorithms} address the second problem by
building the input state from physical principles instead of a generic
circuit.  Two such methods have been proposed recently.
\begin{itemize}
  \item \textbf{Resolution refinement}~[6].  An eigenstate of a
  low-resolution Hamiltonian, defined in a small model space or on a
  coarse lattice, is prepared with any method.  It is then lifted to a
  higher resolution and evolved adiabatically into the corresponding
  eigenstate of the more accurate Hamiltonian.
  \item \textbf{Nucleation}~[7].  The exact ground state of a small
  lattice Hamiltonian is prepared in a circuit, and the lattice is
  then grown to its final size in a series of stages, each carried out
  by Trotterized adiabatic evolution.
\end{itemize}
Their use as input to the rodeo algorithm has not been explored in
these works.

This master's thesis aims to (i) reproduce and critically evaluate the
published one-qubit results, (ii) extend the RA to multi-qubit
Hamiltonians of interacting fermions, (iii) determine how the input
state, including physics-based input states, controls the success
probability and the total cost of the algorithm, (iv) compare the RA
with alternative quantum algorithms, and (v) investigate its practical
performance on modern quantum hardware with noise-mitigation
techniques.  Parts (i) and (ii) are common to all versions of the
project.  After these, the emphasis can be put on hardware execution
and error mitigation, on the input states, or on the comparison of
algorithms, depending on the interests of the student and on the
access to quantum devices.  The thesis involves substantial
computational work, both in classical simulations and in executions on
quantum computers.

\subsection*{Models and Methods}

\textbf{The rodeo algorithm.}  In cycle $n$ an auxiliary qubit
controls the time evolution of the system under $H-E$ for a time
$t_n$ and is then measured.  A new auxiliary qubit is used in each
cycle, or a single one is measured and reset.  For an input state
$\sum_jc_j|E_j\rangle$ the probability that all $N$ cycles succeed is
\[
  P_N(E)=\sum_j|c_j|^2\prod_{n=1}^{N}\cos^2\!\Bigl[\tfrac12(E_j-E)\,t_n\Bigr].
\]
With times drawn from a Gaussian distribution of width $\sigma$, the
average over the times is
\[
  \overline{P_N}(E)=\sum_j|c_j|^2
  \left[\frac{1+e^{-(E_j-E)^2\sigma^2/2}}{2}\right]^{N}.
\]
A scan over the target energy $E$ therefore shows a peak at every
eigenvalue, with a height given by the weight $|c_j|^2$ and a width of
order $1/(\sigma\sqrt N)$.  Components away from the target energy are
suppressed exponentially in $N$, while the weight $|c_0|^2$ of the
wanted state sets the overall success probability.  Random times give
large fluctuations in the suppression, and fixed sequences of times
perform better~[5].

\textbf{Expectation values.}  The expectation value of an operator $O$
in an eigenstate follows from the Hellmann--Feynman theorem,
$\langle E_j|O|E_j\rangle=dE_j(\phi)/d\phi$ at $\phi=0$, where
$E_j(\phi)$ is the eigenvalue of $H+\phi\,O$ determined with the
RA~[3].

\textbf{Model Hamiltonians.}  The project starts from the one-qubit
and two-qubit Hamiltonians of Refs.~[3,4] and then turns to two models
of interacting fermions.  \emph{The pairing Hamiltonian} is a minimal
model of nuclear pairing and superconductivity,
\[
  H=\sum_{p}\varepsilon_p N_p-G\sum_{p,q}P_p^\dagger P_q ,
\]
with single-particle energies $\varepsilon_p$, number operators $N_p$,
pair creation operators $P_p^\dagger$ and pairing strength $G$.
Without broken pairs one qubit per level suffices, and the interaction
becomes $-\tfrac{G}{2}\sum_{p<q}(X_pX_q+Y_pY_q)$ plus one-qubit terms.
The model is exactly solvable~[8,9].  \emph{The Hubbard model in one
dimension},
\[
  H=-t\sum_{i,\sigma}\bigl(c^\dagger_{i\sigma}c_{i+1\sigma}+\mathrm{h.c.}\bigr)
  +U\sum_i n_{i\uparrow}n_{i\downarrow},
\]
needs two qubits per site after a Jordan--Wigner transformation and
is solvable by the Bethe ansatz~[10].  Two to four pairing levels and
a two-site Hubbard model use two to four system qubits and are within
reach of present hardware.  Larger systems are studied in simulation.

\textbf{Input states and cost.}  The input states to be compared are
product and Hartree--Fock states, variational states, states from a
direct adiabatic evolution, and the states from resolution refinement
and nucleation.  The figure of merit is the total cost of obtaining an
eigenstate with a given fidelity: the gates used for the input state
and for the rodeo cycles, divided by the success probability.

\subsection*{Project Outline}

\begin{enumerate}

\item \textbf{Theoretical Foundations}
  \begin{itemize}
    \item Study the rodeo algorithm, its stochastic structure, controlled time evolution, and peak-finding strategy.
    \item Review competing algorithms: quantum phase estimation (QPE) and iterative QPE~[2,11], the variational quantum eigensolver (VQE)~[12], and quantum Krylov methods~[13].
    \item Analyze the Hellmann--Feynman theorem in the context of quantum computing.
    \item Study adiabatic state preparation and the adiabatic theorem~[14], Trotter errors~[15], and the two physics-based methods of Refs.~[6,7].
    \item Review the pairing Hamiltonian and the Hubbard model, their exact solutions and their mapping to qubits; background material is found in the lecture notes of Ref.~[17].
  \end{itemize}

\item \textbf{Reproducing the One-Qubit Demonstration}
  \begin{itemize}
    \item Implement the RA for general one-qubit Hamiltonians.
    \item Classically simulate the success-probability distribution $P_N(E)$ using randomized time evolutions.
    \item Run the circuits on IBM Quantum backends and compare the raw data with the published results~[3].
    \item Integrate measurement error mitigation and study its effects.
    \item Compute expectation values with the Hellmann--Feynman approach and compare with direct preparation of the eigenvector.
  \end{itemize}

\item \textbf{Extension to Multi-Qubit Hamiltonians}
  \begin{itemize}
    \item Reproduce the two-qubit calculation of Ref.~[4] and implement controlled reversal gates.
    \item Construct the pairing Hamiltonian with two to four levels and the two-site Hubbard model, with their exact solutions as reference.
    \item Implement Trotterized time evolution and analyze the Trotter errors.
    \item Evaluate the ability of the RA to resolve multi-level spectra, including nearly degenerate levels, as a function of the pairing strength $G$ and of $U/t$.
    \item Explore the scaling of the RA performance with system size in state-vector simulations.
  \end{itemize}

\item \textbf{Input States for the Rodeo Algorithm}
  \begin{itemize}
    \item Determine the overlap with the exact ground state and the preparation cost for product and Hartree--Fock states, variational states, and direct adiabatic evolution from a simple Hamiltonian.
    \item Implement resolution refinement and nucleation for small cases, following Refs.~[6,7]: add levels around the Fermi level for the pairing Hamiltonian, and grow or refine the chain for the Hubbard model.
    \item Feed the prepared states into the RA and determine success probabilities, the number of cycles needed and the total cost.
    \item Determine how good the input state needs to be: balance the adiabatic evolution time and the number of Trotter steps of the preparation against the number of rodeo cycles.
    \item Extend to excited states, where the input state supplies the overlap and the RA selects the energy.
  \end{itemize}

\item \textbf{Comparison with Other Quantum Algorithms}
  \begin{itemize}
    \item Implement VQE and Krylov-subspace eigensolvers.
    \item Benchmark accuracy, resource cost, and robustness to noise on identical Hamiltonians and with identical input states.
    \item Investigate the operational advantage of the RA under realistic noise.
  \end{itemize}

\item \textbf{Applications on Real Quantum Computers}
  \begin{itemize}
    \item Deploy the algorithms on IBM Quantum and, if possible, on Quantinuum, Amazon Braket or EuroHPC backends.
    \item Systematically compare results: noise, gate depth, mid-circuit measurement quality.
    \item Study the effect of advanced error-mitigation techniques~[16]:
      \begin{itemize}
        \item measurement calibration,
        \item readout symmetrization,
        \item zero-noise extrapolation (ZNE),
        \item probabilistic error cancellation (PEC).
      \end{itemize}
    \item Extract eigenvalues and use the Hellmann--Feynman approach to compute expectation values.
    \item Run small instances (a few levels or sites) of the complete procedure, input-state preparation followed by the RA, and compare the hardware cost of the different input states.
  \end{itemize}

\item \textbf{Exploratory Directions (Optional)}
  \begin{itemize}
    \item Optimization of the RA time-distribution parameters (non-Gaussian choices, fixed and geometric sequences of times~[5]).
    \item Machine-learning tuning of initial states to maximize overlap.
    \item Application to effective Hamiltonians arising from nuclear many-body physics, for example pairing with single-particle energies from the nuclear shell model.
    \item Optimized nucleation~[7], or a combination of nucleation and refinement, as input to the RA.
  \end{itemize}

\end{enumerate}

\subsection*{Milestones and Timeline (12 Months)}

\begin{itemize}
  \item \textbf{Month 1--2:} Literature review; basic RA implementation in classical simulation; reproduction of theoretical formulas.
  \item \textbf{Month 3--4:} Reproduction of the one-qubit energy-scan experiment using simulators; first runs on IBM Quantum; Hellmann--Feynman calculations; error-mitigation benchmarking; detailed comparison with the article data.
  \item \textbf{Month 5--6:} Two-qubit Hamiltonian with controlled reversal gates; pairing and Hubbard Hamiltonians with two to four qubits; Trotterization; classical and quantum benchmarks.
  \item \textbf{Month 7--8:} Input states: standard states, resolution refinement and nucleation for small systems; overlaps, total cost and scaling with system size.
  \item \textbf{Month 9:} Comparison with VQE, Krylov, and QPE approaches on identical Hamiltonians.
  \item \textbf{Month 10--11:} Hardware-based performance study of the multi-qubit cases; statistical scaling; noise analysis; one exploratory direction depending on progress.
  \item \textbf{Month 12:} Final analysis; writing and preparation of the thesis.
\end{itemize}
The weight given to months 7--11 depends on the emphasis chosen after
the common part.

\subsection*{Expected Outcomes}

\begin{itemize}
  \item A full computational implementation of the rodeo algorithm.
  \item Validation and extension of published one-qubit results.
  \item An exploration of the RA applied to multi-qubit Hamiltonians of interacting fermions, in simulation and on real hardware.
  \item A quantitative assessment of how the input state, including states from resolution refinement and nucleation, determines the success probability and the total cost.
  \item Comparative study with VQE, Krylov, and QPE methods under noise.
  \item Assessment of error-mitigation strategies for eigenvalue problems.
  \item A master thesis and, depending on the results, a contribution to a scientific article.
\end{itemize}

\subsection*{References}

\begin{enumerate}
    \item K.~Choi, D.~Lee, J.~Bonitati, Z.~Qian and J.~Watkins, \textit{Rodeo Algorithm for Quantum Computing}, Phys.\ Rev.\ Lett.\ \textbf{127}, 040505 (2021).
    \item D.~Lee, \textit{Quantum techniques for eigenvalue problems}, Eur.\ Phys.\ J.\ A \textbf{59}, 275 (2023), arXiv:2307.03889.
    \item Z.~Qian, J.~Watkins, G.~Given, J.~Bonitati, K.~Choi and D.~Lee, \textit{Demonstration of the rodeo algorithm on a quantum computer}, Eur.\ Phys.\ J.\ A \textbf{60}, 151 (2024).
    \item M.~Bee-Lindgren, Z.~Qian, M.~DeCross, N.~C.~Brown, C.~N.~Gilbreth, J.~Watkins, X.~Zhang and D.~Lee, \textit{Controlled gate networks: theory and application to eigenvalue estimation}, Eur.\ Phys.\ J.\ A \textbf{61} (2025), arXiv:2208.13557.
    \item T.~D.~Cohen and H.~Oh, \textit{Optimizing the rodeo projection algorithm}, Phys.\ Rev.\ A \textbf{108}, 032422 (2023); M.~Patkowski, O.~Ayyildiz, K.~Hunt, N.~Jansen and D.~Lee, \textit{Improved Rodeo Algorithm Performance for Spectral Functions and State Preparation}, arXiv:2602.05978 (2026).
    \item S.~Bogner, H.~Hergert, M.~Hjorth-Jensen, R.~LaRose, D.~Lee and M.~Patkowski, \textit{Quantum State Preparation with Resolution Refinement}, arXiv:2511.14732; Phys.\ Lett.\ B (2026), doi:10.1016/j.physletb.2026.140363.
    \item J.~P.~Sadia, M.~Hjorth-Jensen, D.~Lee and R.~LaRose, \textit{Digital quantum state preparation by nucleation}, arXiv:2609.37527 (2026).
    \item R.~W.~Richardson, \textit{A restricted class of exact eigenstates of the pairing-force Hamiltonian}, Phys.\ Lett.\ \textbf{3}, 277 (1963); J.~Dukelsky, S.~Pittel and G.~Sierra, \textit{Colloquium: Exactly solvable Richardson--Gaudin models for many-body quantum systems}, Rev.\ Mod.\ Phys.\ \textbf{76}, 643 (2004).
    \item D.~J.~Dean and M.~Hjorth-Jensen, \textit{Pairing in nuclear systems: from neutron stars to finite nuclei}, Rev.\ Mod.\ Phys.\ \textbf{75}, 607 (2003).
    \item E.~H.~Lieb and F.~Y.~Wu, \textit{Absence of Mott Transition in an Exact Solution of the Short-Range, One-Band Model in One Dimension}, Phys.\ Rev.\ Lett.\ \textbf{20}, 1445 (1968); F.~H.~L.~Essler \textit{et al.}, \textit{The One-Dimensional Hubbard Model}, Cambridge University Press (2005).
    \item A.~Yu.~Kitaev, \textit{Quantum measurements and the Abelian Stabilizer Problem}, arXiv:quant-ph/9511026 (1995); D.~S.~Abrams and S.~Lloyd, \textit{Quantum Algorithm Providing Exponential Speed Increase for Finding Eigenvalues and Eigenvectors}, Phys.\ Rev.\ Lett.\ \textbf{83}, 5162 (1999).
    \item A.~Peruzzo \textit{et al.}, \textit{A variational eigenvalue solver on a photonic quantum processor}, Nat.\ Commun.\ \textbf{5}, 4213 (2014).
    \item N.~H.~Stair, R.~Huang and F.~A.~Evangelista, \textit{A Multireference Quantum Krylov Algorithm for Strongly Correlated Electrons}, J.\ Chem.\ Theory Comput.\ \textbf{16}, 2236 (2020).
    \item T.~Albash and D.~A.~Lidar, \textit{Adiabatic quantum computation}, Rev.\ Mod.\ Phys.\ \textbf{90}, 015002 (2018).
    \item A.~M.~Childs, Y.~Su, M.~C.~Tran, N.~Wiebe and S.~Zhu, \textit{Theory of Trotter Error with Commutator Scaling}, Phys.\ Rev.\ X \textbf{11}, 011020 (2021).
    \item K.~Temme, S.~Bravyi and J.~M.~Gambetta, \textit{Error Mitigation for Short-Depth Quantum Circuits}, Phys.\ Rev.\ Lett.\ \textbf{119}, 180509 (2017); Z.~Cai \textit{et al.}, \textit{Quantum error mitigation}, Rev.\ Mod.\ Phys.\ \textbf{95}, 045005 (2023).
    \item M.~Hjorth-Jensen, \textit{Quantum Computing and Quantum Machine Learning}, lecture notes and codes,\\ \url{https://github.com/CompPhysics/QuantumComputingMachineLearning}.
\end{enumerate}

\end{document}
